The next quotient has an unconditional isomorphism-lifting property.
Lifting objects requires an additional hypothesis.

\begin{theorem}[a compatible Jacobson quotient]
\label{thm:upper-jacobson-quotient}
For any unital inverse semiring $T=(Q,F)$, define
\[
 x\in C_q\quad\Longleftrightarrow\quad
 t_{pqr,e}(uxv)\in\operatorname{Jac}(F(e))
\]
for every $u\in F(p)$, $v\in F(r)$ and $e\in U(Q)$ with
$pqr\leq e$. Then $C$ is the largest compatible coefficient ideal
with $C_e\subseteq\operatorname{Jac}(F(e))$ at every upper
idempotent. It contains $H$, and $T/C$ has no nonzero compatible
coefficient ideal with this property. It commutes entrywise with
every finite matrix construction.

For each finite idempotent support matrix $E$, the kernel $N_E$ of
$A_E(T)\to A_E(T/C)$ is a Jacobson radical, possibly nonunital,
ring. For any idempotent $P\in A_E(T)$,
\[
 P N_E P\subseteq\operatorname{Jac}(P A_E(T)P).
\]
Every isomorphism between reductions of two finitely generated
projective $T$-semimodules lifts to an isomorphism. Thus
$V(T)\to V(T/C)$ is injective, and a quotient projective has at
most one lifting isomorphism class.

If all $F(e)$, $e\in U(Q)$, are semiperfect, every quotient
projector lifts at its exact support. In that case
\[
 V(T)\simeq V(T/C),\qquad K_0(T)\simeq K_0(T/C).
\]
There is no finiteness or commutativity hypothesis on $Q$ or $F$.
\end{theorem}

The unrestricted assertion gives neither existence of all lifting
objects nor an injection on group completions. The ordinary
semilocal domain in Example~\ref{ex:jacobson-lifting-boundary}
already prevents unconditional object lifting. This is a quotient
of coefficient groups retaining $Q$; it is not the Bourne
quotient by either of the usual global semiring Jacobson radicals.


Corollary~\ref{cor:jacobson-transition-tests} evaluates this same
coefficient ideal by upward ordinary-ring transitions and explicit
two-by-two fiber rings, with finitely many tests when $Q$ is finite.
